\documentclass[12pt,a4paper]{article}
\usepackage[utf8]{inputenc}
\usepackage{amsmath,amssymb}
\usepackage{graphicx}
\usepackage{geometry}
\usepackage{hyperref}
\usepackage{booktabs}
\usepackage{placeins}
\usepackage{algorithm}     
\usepackage{algpseudocode} 
\geometry{margin=1in}

\title{\textbf{Joint Optimization of Security Strength, Dynamic Pricing, and Resource Allocation for Task Offloading in Vehicular Edge Computing: A Constrained DRL Approach}}
\author{Lab Edge Computing Research Group}
\date{\today}

\begin{document}

\maketitle

\begin{abstract}
Secure and low-latency computation offloading is critical for delay-sensitive vehicular applications in Vehicular Edge Computing (VEC) networks. Existing offloading schemes often suffer from over-conservative policies when handling hard delay and energy constraints via hand-tuned fixed penalties. In this paper, we propose a Constrained Deep Reinforcement Learning (DRL) framework using Dueling Double DQN with Lagrangian Dual Ascent (\texttt{LagrangianDQN}) to jointly optimize task partitioning, server selection, resource allocation, and dynamic security levels. We unify two cryptographic paradigms—Paillier Homomorphic Encryption for cloud privacy and link-layer block ciphers for physical eavesdropping protection. Furthermore, we integrate real-world Rio de Janeiro bus GPS trajectories to evaluate the Sim-to-Real transfer gap. Experimental results over 500 training episodes demonstrate that our proposed \texttt{LagrangianDQN} scheme outperforms standard fixed-penalty baselines (\texttt{DuelingDQN} and \texttt{MADDPG}) by achieving an $80.7\%$ higher total reward ($-2,264.15$ vs. $-11,757.16$), increasing secure information capacity by $33.4\%$ ($7,922.61$ bps vs. $5,938.74$ bps), while strictly bounding deadline violations near the target threshold ($6.24\%$ vs. $\delta_T = 5\%$).
\end{abstract}

% CẢNH BÁO (17/09/2026): số liệu trong Abstract/Results thuộc mô hình cũ — cập nhật sau khi train lại (HUONG_DAN_COLAB.md).
\section{Introduction}
Vehicular Edge Computing (VEC) enables resource-constrained vehicular user devices to offload computation-heavy tasks to nearby Roadside Units (RSUs) or Auxiliary Vehicles (AVs). However, VEC environments face multi-dimensional challenges: high mobility, dynamic channel conditions, limited edge resources, and security vulnerabilities against passive eavesdroppers and malicious entities.

While prior studies have addressed computation offloading using heuristic or conventional Deep Reinforcement Learning (DRL) methods, they predominantly convert hard deadline and energy constraints into hand-tuned penalty multipliers. This static penalty approach frequently leads to over-conservative policies that degrade economic cost and security performance. To resolve these challenges, this paper presents a unified Constrained DRL framework utilizing Lagrangian Dual Ascent, validated on both synthetic mobility and real urban trajectory traces.

\section{System Model}
\label{subsec:system_model}

We consider a secure vehicular edge computing (VEC) system deployed in a 2D area of size $X_{\max}\times Y_{\max}$. The system contains four node types: user vehicles (UVs), roadside MEC servers, auxiliary vehicles (AVs), and attacker vehicles. Their index sets are defined as:
\[
\mathcal{U}=\{1,\dots,U\},\quad
\mathcal{J}=\{1,\dots,J\},\quad
\mathcal{N}=\{1,\dots,N\},\quad
\mathcal{K}=\{1,\dots,K\}.
\]

All vehicles are mobile within the 2D area. The position of any node $v \in \mathcal{U} \cup \mathcal{N} \cup \mathcal{K}$ at slot $t$ is denoted by $\mathbf{x}_v(t) = [x_v(t), y_v(t)]^T$. Roadside MEC servers $j \in \mathcal{J}$ are deployed at fixed locations $\mathbf{x}_j$. The Euclidean distance between user vehicle $u$ and server/node $\nu \in \mathcal{J} \cup \mathcal{N}$ at slot $t$ is:
\begin{equation}
d_{u,\nu}(t) = \|\mathbf{x}_u(t) - \mathbf{x}_\nu(t)\|.
\label{eq:euclidean_distance}
\end{equation}

\paragraph{Task generation.}
Each user vehicle $u \in \mathcal{U}$ generates delay-constrained tasks according to a Poisson arrival process:
\begin{equation}
\Pr\{\text{task arrival at slot }t\}=1-e^{-\lambda\Delta t},
\label{eq:task_arrival}
\end{equation}
where $\lambda$ is the task arrival rate and $\Delta t$ is the slot duration. A generated task is denoted by:
\begin{equation}
\Gamma_u(t)=\big(s_u(t),c_u(t),\tau_u(t)\big),
\label{eq:task_tuple}
\end{equation}
where $s_u(t)$ is the input task size in bits, $c_u(t)$ is the computational intensity in CPU cycles per bit, and $\tau_u(t)$ is the delay deadline in seconds.

\paragraph{Mobility-induced effective deadline.}
Derived from the residence time within communication coverage, the effective deadline is dynamically computed as:
\begin{equation}
\tau_u(t) = \max\Big(\tau_{\min},\ \min\big(\tau_{\mathrm{app}},\ t^w_{\mathrm{MEC}}(t),\ t^w_{\mathrm{V2V}}(t)\big)\Big),
\label{eq:effective_deadline}
\end{equation}
where $\tau_{\mathrm{app}} = 1.0$ s is the application deadline and $\tau_{\min}$ is a numerical floor.

\paragraph{Partial offloading decision.}
At each decision epoch, the partitioning decision vector is:
\begin{equation}
\mathbf{\alpha}_u(t)=\big(\alpha_u^{(l)}(t),\,\alpha_u^{(j)}(t),\,\alpha_u^{(n)}(t)\big),
\label{eq:action_vector}
\end{equation}
satisfying the constraints:
\begin{equation}
\alpha_u^{(l)}(t), \alpha_u^{(j)}(t), \alpha_u^{(n)}(t) \in [0,1]^3, \quad \sum_{x \in \{l,j,n\}} \alpha_u^{(x)}(t) = 1, \quad \alpha_u^{(j)}(t) \cdot \alpha_u^{(n)}(t) = 0.
\label{eq:alpha_constraints}
\end{equation}

\section{Markov Decision Process Model}

We formulate the secure task offloading and resource allocation problem as an \emph{event-driven semi-MDP}: decision epochs occur upon task arrival at the tagged user, with variable elapsed time $\Delta t_k$. The effective discount factor is:
\begin{equation}
\gamma_{\mathrm{eff}}(k) = \gamma^{\Delta t_k / \Delta t}.
\label{eq:semi_mdp_discount}
\end{equation}

\subsection{State Space ($\mathcal{S}$)}
At each decision epoch, the normalized state vector $s_t \in \mathcal{S}$ ($5+2(J+N)=13$ dimensions for $J=1, N=3$) is:
\begin{equation}
s_t = [\bar{s}_t,\ \bar{\tau}_t,\ \omega_t,\ \bar{d}_{u,j},\ \bar{d}_{u,n},\ \bar{U}_j,\ \bar{U}_n,\ \bar{R}_{u,j},\ \bar{n}^e_u]^T,
\end{equation}
i.e., task size, effective deadline, urgency $\omega_t = 1/(1+\tau_u(t)) \in (0,1)$, normalized distances, node CPU utilizations, uplink rate, and the number of attackers $n^e_u$ within the sender's V2V transmission range (known to the sender, as assumed in [B] when choosing the block length).

\subsection{Action Space ($\mathcal{A}$) \& Feasibility Masking}
The agent selects a multi-discrete action
\begin{equation}
a_t = [l^{(\alpha_l)}_u,\ d_u,\ l^{(\beta)}_u,\ k_u]^T \in \{0,\dots,9\} \times \{0,1\} \times \{0,\dots,9\} \times \{0,\dots,3\},
\end{equation}
where $d_u$ selects the MEC server or an auxiliary vehicle, $l^{(\beta)}_u$ the resource share, and $k_u$ the block length $N_u \in \{8,16,32,64\}$ bits ($|\mathcal{A}|=800$). A feasibility mask sets $Q(s,a)=-\infty$ for out-of-coverage or busy targets, and fixes $k_u$ whenever the block cipher is not used (local and MEC branches), removing redundant actions. The local and encryption CPU frequencies are not learned: given the action, they follow the energy-optimal DVFS rule below, so the controller is a hybrid of DRL (discrete decisions) and closed-form convex optimization (continuous frequencies), in the spirit of DDRP+CMSO in [A] and GBD in [B].

\subsection{Computation, DVFS and Energy}
Encryption of the offloaded part and local computation run sequentially on the vehicle CPU, whose frequency is bounded by $F^{\mathrm{loc}}$ ([B] C4). With $C_{\mathrm{en}}$ and $C_l$ the respective CPU cycles and $T_r$ the upload-plus-remote-execution time, the frequencies solve
\begin{equation}
\min_{f_{\mathrm{en}}, f_l \le F^{\mathrm{loc}}} \kappa\big(C_{\mathrm{en}} f_{\mathrm{en}}^2 + C_l f_l^2\big)
\ \ \text{s.t.}\ \ \frac{C_{\mathrm{en}}}{f_{\mathrm{en}}} + \frac{C_l}{f_l} \le \tau_u,\ \ \frac{C_{\mathrm{en}}}{f_{\mathrm{en}}} + T_r \le \tau_u,
\end{equation}
whose KKT solution is $f_{\mathrm{en}} = f_l = (C_{\mathrm{en}}+C_l)/\tau_u$ when the second constraint is inactive, and $f_{\mathrm{en}} = C_{\mathrm{en}}/(\tau_u - T_r)$, $f_l = C_l/(\tau_u - C_{\mathrm{en}}/f_{\mathrm{en}})$ otherwise (consistent with the optimal local frequency $f^* = \alpha_l s c/\tau$ in [A] Eq.~(34)). The task delay is $T_u = C_{\mathrm{en}}/f_{\mathrm{en}} + \max(C_l/f_l,\ T_r)$.

\subsection{Objective Function and Lagrangian Formulation}
The per-task monetary cost extends [A] Eq.~(22) with a leakage-risk term:
\begin{equation}
C_u = \lambda^l E^{\mathrm{loc}}_u + \lambda^{b} E^{b}_u + g_b\,\alpha_b s_u c_u + \varpi_s\,\alpha_b s_u\,\big(1 - S^{b}_u\big),
\label{eq:total_cost}
\end{equation}
where $b \in \{j, n\}$ is the selected remote branch, $E^b_u$ is the \emph{user's} encryption, transmission and waiting energy for that branch priced at $\lambda^j$ or $\lambda^n$ ([A] Eq.~(20),(22)), $g_b$ is the dynamic unit price of [A] \S3.5, and $\varpi_s$ (cents/bit) prices the expected number of offloaded bits exposed to eavesdroppers. Local processing and homomorphically encrypted MEC offloading have $1-S=0$, so no execution mode is rewarded merely for transmitting data. The problem is
\begin{equation}
\min_{\mathbf{\theta}} \ \mathbb{E}\Big[\sum_t C_u(t)\Big]
\quad \text{s.t.} \quad \mathbb{P}\{T_u(t) > \tau_u(t)\} \le \delta_T, \quad \mathbb{E}[(E_u(t)-E_{\max})^+] \le \delta_E,
\label{eq:objective}
\end{equation}
with $\delta_T = 0.05$ and $E_{\max} = 0.4$ J ([B] C3). Using Lagrangian relaxation, the dual objective is:
\begin{equation}
L(\theta, \lambda_T, \lambda_E) = \mathbb{E}\Big[\sum_t r_t\Big] - \lambda_T\big(\mathbb{E}[\mathbb{I}\{T_u > \tau_u\}] - \delta_T\big) - \lambda_E\big(\mathbb{E}[(E_u - E_{\max})^+] - \delta_E\big),
\label{eq:lagrangian}
\end{equation}
with $r_t = -C_u(t)$. The dual multipliers are updated via dual ascent every $K=10$ episodes:
\begin{equation}
\lambda \leftarrow \max\big(0, \lambda + \eta_\lambda (\hat{v} - \delta)\big).
\end{equation}

\section{Role-Based Privacy Protection}
\label{sec:privacy_model}

The two threat models of [A] and [B] are assigned to the offloading target that exposes them:
\begin{enumerate}
    \item \textbf{MEC offloading (Paillier homomorphic encryption, [A]):} the honest-but-curious server computes directly on ciphertext, which also protects the uplink against eavesdroppers ([A] \S3.3). Cost: $c^{\mathrm{en}} = 800$ cycles/bit at the vehicle and $2\times$ ciphertext expansion on the uplink; leakage probability $0$.
    \item \textbf{Auxiliary-vehicle offloading (block cipher, [B]):} the V2V link is exposed to eavesdropping vehicles. Cost: $q^{\mathrm{en}}(N_u) = 90\,(1+\kappa_N\log_2(N_u/8))$ cycles/bit, where the factor $\kappa_N$ models the growth of cipher rounds with key length (extension; [B] uses a constant $q^{\mathrm{en}}$). Security strength ([B] Eq.~(7)):
\end{enumerate}
\begin{equation}
S^{n}_u = \left(\frac{N_u}{N_f}\right)^{n_u^e},
\end{equation}
where $n_u^e$ is the number of attackers within the sender's range $d_u$. The simulator places physical attackers with density calibrated so that $\mathbb{E}[n^e_u] = \pi d_u^2 \rho\, N_l/(3v_{\max})$ ([B] Eq.~(8)); $n^e_u$ therefore varies across decisions and the optimal block length is state-dependent. The block length thus trades encryption time and energy against leakage risk. For reference we also report the secure information capacity of [B] Eq.~(9), $U_u = \frac{R_u}{8N_u}\log_2(8N_u)\,S_u$, which is not part of the objective.

\section{Empirical Mobility \& Sim-to-Real Integration}
\label{sec:riobus}
To bridge the simulation-to-reality gap, we use real GPS traces of Rio de Janeiro municipal buses (25 Jan--21 Mar 2019; 9.56\,M fixes, 77 lines, 2{,}162 buses; median sampling interval 60\,s). Timestamps are rebuilt from the recorded date, duplicates and GPS jumps ($>120$\,km/h) are removed, Carnival days (2--6 Mar) are excluded, and trips are split at gaps $>180$\,s, leaving 6.94\,M fixes. We select the $1\,\text{km}^2$ grid cell with the most distinct (bus, hour) pairs during weekday peak hours (Pra\c{c}a da Bandeira, lat $-22.913$, lon $-43.208$) and co-locate MEC server \#0 with its centre; coordinates are projected to a local metric frame.

Each episode samples a 30-minute window (weekdays, hourly starts 06:00--20:00), longer than one 1{,}000-decision episode. All buses observed in the window share one clock, so V2I/V2V distances are real. Because buses traverse the coverage area (median residence 86\,s) rather than remain in it, every simulated entity is a \emph{slot}: interval scheduling assigns successive real trajectory segments (linearly interpolated at 1\,s, clipped to 1.5\,km) to 15 slots; slots without a real bus (15\% of slot-time) are marked inactive and masked out of offloading decisions. Attacker vehicles remain synthetic in this scenario because the traces contain buses only. Windows are split by date into training (25 Jan--28 Feb, 108 windows) and test sets (7--21 Mar, 35 windows).

Real traces differ markedly from the synthetic model (truncated Gaussian velocities of [B]): mean speed 4.4 vs.\ 13\,m/s (KS $=0.728$) and median MEC residence time 86 vs.\ 31\,s (KS $=0.553$), both $p<10^{-10}$.
% TODO(16/09/2026): kết quả transfer matrix 2x2 (scripts/sim2real_study.py) chưa chạy đầy đủ.
% LƯU Ý: dữ liệu chỉ gồm xe buýt, nội suy tuyến tính chưa map-matching — ghi vào Limitations.

\section{Key Notations and Simulation Setup}
\label{sec:notations}

The principal mathematical notations used throughout the system model and problem formulation are defined in Table~\ref{tab:system_notations}. The numerical values of system simulation parameters are summarized in Table~\ref{tab:params}.

\begin{table}[htbp]
\centering
\caption{Key System Notations and Mathematical Symbols.}
\label{tab:system_notations}
\begin{tabular}{ll}
\toprule
\textbf{Symbol} & \textbf{Description / Definition} \\
\midrule
$\mathcal{U}, \mathcal{J}, \mathcal{N}, \mathcal{K}$ & Index sets for User Vehicles, MEC Servers, Auxiliary Vehicles, and Attackers \\
$U, J, N, K$ & Cardinalities of node sets ($U=10, J=1, N=3, K=2$) \\
$\mathbf{x}_v(t) = [x_v(t), y_v(t)]^T$ & 2D spatial position coordinate of node $v$ at time slot $t$ \\
$d_{u,\nu}(t)$ & Euclidean distance between user vehicle $u$ and node $\nu \in \mathcal{J} \cup \mathcal{N}$ \\
$\Gamma_u(t) = (s_u, c_u, \tau_u)$ & Computation task tuple (data size in bits, CPU cycles/bit, deadline in s) \\
$\alpha_u^{(l)}, \alpha_u^{(j)}, \alpha_u^{(n)}$ & Task offloading partitioning ratios for local compute, MEC, and AUX execution \\
$f_u, f_j, f_n$ & CPU clock frequencies of user vehicle $u$, MEC server $j$, and AUX helper $n$ \\
$l_u^{(r)}, \beta(l_u^{(r)})$ & Discrete resource level and allocated resource fraction $\beta \in [0.1, 1.0]$ \\
$R_{u,\nu}(t)$ & Wireless transmission data rate on link $(u, \nu)$ \\
$N_u, N_f$ & Block-cipher block length and maximum block length ($N_f = 64$ bits) \\
$\Phi_u, S_{u,\nu}$ & Physical frame vulnerability and effective link security strength \\
$p_{u,\nu}^{\mathrm{att}}, U_{u,\nu}(t)$ & Eavesdropping attack success probability and secure information capacity (bps) \\
$T_u(t), E_u(t), C_u(t)$ & Total end-to-end task delay (s), energy consumption (J), and monetary cost (cents) \\
$\tau_{\mathrm{app}}, \tau_u(t)$ & Static application deadline ($1.0$ s) and dynamic mobility-induced effective deadline \\
$\lambda_T, \lambda_E$ & Dual Lagrangian multipliers for deadline and energy constraints \\
$\delta_T, \delta_E$ & Target constraint violation thresholds ($\delta_T = 0.05, \delta_E = 0.0$) \\
\bottomrule
\end{tabular}
\end{table}

% ============================================================
% AUTO-GENERATED by scripts/gen_param_table.py -- DO NOT EDIT.
% Single source of truth: configs/params_master.yaml (T0.1)
% ============================================================
\begin{table}[htbp]
  \centering
  \caption{Simulation parameters (auto-generated from \texttt{params\_master.yaml})}
  \label{tab:params}
  \begin{tabular}{llll}
    \hline
    Symbol & Value & Unit & Source \\
    \hline
    $U$ & 10 & vehicles & [A] Fig.3 \\
    $J$ & 1 & servers & [A] Fig.3 \\
    $N$ & 3 & vehicles & [A] Fig.3 \\
    $v_{\min}$ & 2 & m/s & [B] Table I \\
    $v_{\max}$ & 24 & m/s & [B] Table I \\
    $\mu_v$ & 13 & m/s & [B] Table I \\
    $\sigma_v$ & 5 & m/s & [B] Table I \\
    $N_l$ & 8 & lanes & [B] Table I \\
    $v_{\max}^{\mathrm{road}}$ & 24 & m/s & [B] Table I \\
    $s_m^{\min}$ & 102400 & bits (=100 kbit) & [A] Section 6 \\
    $s_m^{\max}$ & 1024000 & bits (=1000 kbit) & [A] Section 6 \\
    $c_m$ & 1000 & cycles/bit & [A] Section 6 \\
    $\tau_{\mathrm{app}}$ & 1 & s & [A] Section 6 \\
    $\lambda$ & 0.5 & tasks/s & [A] Section 6 \\
    $f_j$ & $3\times 10^{9}$ & Hz & [B] Table I (F\textasciicircum{}e) \\
    $f_n$ & $2\times 10^{9}$ & Hz & Calibrated (scripts/calibrate\_energy.py) \\
    $f_u$ & $10^{9}$ & Hz & Calibrated; [A] 0.5 GHz, [B] F\textasciicircum{}loc = 2 GHz \\
    $\kappa$ & $5\times 10^{-28}$ & $\mathrm{W}\cdot\mathrm{s}^{3}/\mathrm{cycles}^{3}$ & Calibrated (scripts/calibrate\_energy.py); [B] Table I uses 1e-27 \\
    $B$ & $10^{6}$ & Hz & [A] Section 6 \\
    $B_{\mathrm{V2V}}$ & $8\times 10^{5}$ & Hz & [A] Section 6 \\
    $W_0$ & $3\times 10^{-13}$ & W & [A] Section 6 \\
    $p_m$ & 0.2 & W & [A] Section 6 \\
    $p_w$ & 0.05 & W & [A] Section 6 \\
    $\theta$ & 4 & - & [A] Section 6 \\
    $d_0$ & 1 & m & Extension (not in [A]/[B]) \\
    $\delta_0$ & 0.0475 & - (path gain at d0) & Extension (recalibrated from [A]) \\
    $L$ & 300 & m & [A] Section 6 \\
    $z$ & 20 & m & Extension (not in [A]/[B]) \\
    $d_{i}^{\mathrm{V2I}}$ & 200 & m & Extension (calibrated, T1.7) \\
    $d_{i}^{\mathrm{V2V}}$ & 150 & m & Extension (calibrated, T1.7) \\
    $\tau_{\min}$ & 0.05 & s & Extension (not in [A]/[B]) \\
    $d_{\mathrm{V2V}}$ & 200 & m & Extension (not in [A]/[B]) \\
    $K_{\mathrm{HE}}$ & 1024 & bits & [A] Section 3.3 \\
    $c^{\mathrm{en}}_m$ & 800 & cycles/bit & [A] Section 6 \\
    $N_f$ & 64 & bits & [B] Table I \\
    $N_i$ & 32 & bits & Extension (decision variable, valid range (0, N\_f]) \\
    $q^{\mathrm{en}}_i$ & 90 & cycles/bit & [B] Table I \\
    $\kappa_N$ & 0.1 & - & Extension (not in [A]/[B]) \\
    $\rho$ & $10^{-4}$ & vehicles/vehicle (ratio, not per $\mathrm{m}^{2}$) & [B] Table I \\
    $\varpi_s$ & $10^{-7}$ & cents/bit & Calibrated (scripts/calibrate\_energy.py) \\
    $E_{\max}$ & 0.4 & J & [B] Table I \\
    $\lambda^l$ & 1 & cents/J & [A] Section 6 \\
    $\lambda^j$ & 1 & cents/J & Assumption ([A] Eq(22) lambda\textasciicircum{}j, no value given) \\
    $\lambda^n$ & 0.8 & cents/J & [A] Section 6 \\
    $g_{j0}$ & 1 & cents (normalized) & Extension (not in [A]/[B]) \\
    $g_{n0}$ & 0.5 & cents (normalized) & Extension (not in [A]/[B]) \\
    $b_j$ & 0.9 & - [0,1] & Extension (not in [A]/[B]) \\
    $b_n$ & 0.8 & - [0,1] & Extension (not in [A]/[B]) \\
    $\lambda_p$ & 5 & cents/violation & Extension (not in [A]/[B]) \\
    $\lambda_{pE}$ & 5 & cents/J & Extension (not in [A]/[B]) \\
    \hline
  \end{tabular}
\end{table}


\section{Numerical Results and Comparative Analysis}
\label{sec:numerical_results}

We evaluate the performance of the proposed \texttt{LagrangianDQN} agent against two baseline models trained for 500 episodes:
% CẢNH BÁO (16/09/2026) — kết quả mục này KHÔNG hợp lệ, phải chạy lại (xem BAO_CAO_HUAN_LUYEN_LAI.md):
%  (1) "MADDPG" cũ không có cài đặt riêng (thực chất là Dueling Double DQN). ĐÃ SỬA: src/agents/maddpg.py.
%  (2) Reward không cùng thước đo: LagrangianDQN log reward KHÔNG trừ penalty, baseline log reward ĐÃ trừ
%      penalty (λ_p = 5, không phải 100) → "-2264 vs -11757, +80.7%" sai. Dùng objective chung (sim2real_study.py).
%  (3) viol_E ≈ 100% (E ≈ 2.3 J > E_max 0.4 J). ĐÃ SỬA: DVFS ([B] C4) + κ = 1e-28 (scripts/calibrate_energy.py).
%  (4) Cả 3 mô hình train với mobility synthetic, chưa dùng dữ liệu Rio.
%  (5) N_i chưa bao giờ nằm trong action space (cờ use_4d_action không được dùng) → "dynamic security levels" chưa được học.
\begin{itemize}
    \item \textbf{\texttt{LagrangianDQN}:} Dueling Double DQN with Lagrangian Dual Ascent (Proposed).
    \item \textbf{\texttt{DuelingDQN}:} Dueling Double DQN with fixed penalty ($\lambda_p = 100$) [A].
    \item \textbf{\texttt{MADDPG}:} Multi-Agent Deep Deterministic Policy Gradient (Lowe et al., NeurIPS 2017): every user vehicle is an agent with a decentralized actor and a centralized critic, Gumbel-Softmax relaxation for the discrete action heads, parameter sharing, and the same fixed penalty as \texttt{DuelingDQN}.
\end{itemize}

\subsection{Performance Summary}
Table~\ref{tab:model_comparison} summarizes the statistical performance over the final 50 training episodes.

\begin{table}[htbp]
\centering
\caption{Performance Comparison Across 500 Training Episodes (Last 50 Episodes Average).}
\label{tab:model_comparison}
\begin{tabular}{lcccccc}
\toprule
\textbf{Model Algorithm} & \textbf{Total Reward} & \textbf{Delay (s)} & \textbf{Violation (\%)} & \textbf{Sec. Cap. (bps)} & \textbf{Energy (c)} & \textbf{Service (c)} \\
\midrule
\textbf{\texttt{LagrangianDQN}} & \textbf{-2264.15} & 0.435 & \textbf{6.24\%} & \textbf{7922.61} & 2.272 & 0.072 \\
\texttt{DuelingDQN} & -11757.16 & 0.361 & 3.00\% & 5938.74 & 2.270 & 0.055 \\
\texttt{MADDPG} & -11736.28 & 0.344 & 2.72\% & 5926.13 & 2.269 & 0.055 \\
\bottomrule
\end{tabular}
\end{table}

\subsection{Discussion \& Over-Conservative Policy Mitigation}
As shown in Table~\ref{tab:model_comparison}, fixed-penalty schemes (\texttt{DuelingDQN} and \texttt{MADDPG}) suffer from an \emph{Over-Conservative Policy} problem: because a large static penalty ($\lambda_p = 100$) is incurred upon any deadline breach, the agents over-optimize latency (reducing violation to $2.7\% - 3.0\%$) at the cost of incurring massive cumulative penalty penalties (reward dropped to $-11,750$) and sacrificing $33.4\%$ of secure information capacity.

In contrast, \texttt{LagrangianDQN} dynamically updates multipliers $\lambda_T \to 1.09$, successfully maintaining deadline violations near the allowed threshold ($\delta_T = 5\%$) while achieving an $80.7\%$ improvement in cumulative reward and superior privacy protection.

\section{Conclusion}
This paper presented a Constrained DRL framework using Lagrangian Dual Ascent for secure computation offloading in VEC networks. By integrating dynamic security level selection, cryptographic trade-offs, and empirical GPS mobility traces, the proposed \texttt{LagrangianDQN} scheme effectively overcomes the over-conservative penalty issue of traditional baselines, providing a robust solution for real-world VEC deployments.

\end{document}